\documentclass[12pt]{article}
\usepackage[T1]{fontenc}
\usepackage[utf8]{inputenc}
\usepackage{lmodern}
\usepackage{textcomp}
\usepackage{enumitem}
\usepackage{geometry}
\geometry{margin=1in}
\begin{document}
\begin{center}
Answer Key - Business Administration - Undergraduate Level Assessment 10 \
\small (Answers are ordered exactly as in the question paper.)
\end{center}

\section*{Multiple Choice}
\begin{enumerate}[label=\arabic*.]
\item \textbf{Answer:} D
\\ \textit{Options:} A) Good usage.; B) Servicescape.; C) Service mix.; D) Service encounters.
\\ \textit{Note:} Service encounters refer to the specific moments when customers engage with a service provider, making them the most relevant term to describe the direct interactions that occur during the service delivery process.
\item \textbf{Answer:} B
\\ \textit{Options:} A) The firm must avoid focusing on non-variables such as profitability and volume.; B) The market segment must have measurable purchasing power and size.; C) The company must expand beyond its marketing capabilities to capture growing markets.; D) The market segment must reflect the population's changing attitudes and lifestyles.
\\ \textit{Note:} The correct answer emphasizes that for market segmentation to be effective, it is essential to identify segments that can be quantified in terms of their purchasing power and size, ensuring that the business can target and meet the needs of viable consumer groups.
\item \textbf{Answer:} D
\\ \textit{Options:} A) Proximity; B) Timeliness; C) Prominence; D) Impact
\\ \textit{Note:} The criterion of impact is most important in public relations messages because it emphasizes the significance of the information being communicated and its potential effects on the target audience, which is crucial for shaping public perception and influencing behavior.
\item \textbf{Answer:} D
\\ \textit{Options:} A) Social; B) Business; C) Organisational; D) Company
\\ \textit{Note:} The correct answer is "Company" because the layers of cultural influence typically include individual, group, and organizational culture, while "Company" refers to a specific entity rather than a layer of cultural influence.
\item \textbf{Answer:} B
\\ \textit{Options:} A) Telephone; B) Face-to-face; C) E-mail; D) Video conference
\\ \textit{Note:} Face-to-face communication allows for the exchange of non-verbal cues, immediate feedback, and a deeper personal connection, making it the richest form of information in business interactions.
\end{enumerate}

\section*{True / False}
\begin{enumerate}[label=\arabic*.]
\setcounter{enumi}{5}
\item \textbf{Answer:} False
\\ \textit{Options:} A) True; B) False
\\ \textit{Note:} Impulse products are typically unplanned purchases made on the spur of the moment, requiring little to no pre-purchase thought or evaluation.
\item \textbf{Answer:} False
\\ \textit{Options:} A) True; B) False
\\ \textit{Note:} An international supply chain is dynamic and interconnected, emphasizing collaboration, adaptability, and the integration of various processes beyond just linear operations and transactions.
\item \textbf{Answer:} True
\\ \textit{Options:} A) True; B) False
\\ \textit{Note:} The Chase Jones model of issues management emphasizes that during the strategy step, organizations evaluate their capability to effectively address and lead specific issues, determining their readiness and approach to manage potential challenges.
\item \textbf{Answer:} True
\\ \textit{Options:} A) True; B) False
\\ \textit{Note:} The Sarbanes-Oxley Act establishes strict regulations and accountability measures for financial reporting and corporate governance, thereby promoting transparency and ethical behavior in business practices.
\item \textbf{Answer:} False
\\ \textit{Options:} A) True; B) False
\\ \textit{Note:} Time utility encompasses both the convenience of purchasing products at the right time and the availability of those products at the right location, making the statement misleading by omitting location as a factor in time utility.
\end{enumerate}

\section*{Long Form Response}
\begin{enumerate}[label=\arabic*.]
\setcounter{enumi}{10}
\item \textbf{Answer:} Continuous evaluation of current and future trends is crucial for public relations counselors as it enables them to interpret trends that inform organizational management decisions. By understanding evolving public sentiments, technological advancements, and market dynamics, PR counselors can provide strategic insights that shape an organization's response to challenges and opportunities. This proactive approach not only enhances an organization's reputation but also aligns its goals with stakeholder expectations, ultimately fostering better decision-making. Consequently, organizations that leverage these insights can adapt more effectively to changes in the business environment, ensuring sustained success and relevance.
\\ \textit{Note:} The answer correctly emphasizes that continuous trend evaluation allows public relations counselors to provide strategic insights that guide organizational management decisions, ultimately enhancing reputation and aligning goals with stakeholder expectations.
\item \textbf{Answer:} The Vatican's response to child abuse cover-up accusations in 2010 highlighted a significant tension between accountability and the perception of criticism as an attack on the Catholic Church. Many church officials framed the allegations as assaults on the institution's integrity, which complicated the dialogue around necessary reforms and transparency. This defensive stance not only hindered the Church's ability to address the abuse crisis effectively but also alienated victims and advocates seeking justice. As a result, the Vatican's approach risked reinforcing a culture of silence and avoidance, ultimately affecting the Church's reputation and its relationship with the broader community. This situation illustrates how organizational responses to criticism can impact stakeholder trust and the institution's long-term viability.
\\ \textit{Note:} The answer correctly identifies that the Vatican's defensive framing of child abuse cover-up accusations as attacks on the Church impeded accountability and reform efforts, thereby alienating victims and advocates and complicating the necessary dialogue on addressing the crisis.
\end{enumerate}

\vfill
\noindent\textit{End of Answer Key}
\end{document}